\section{相关工作}

\paragraph{在线话语中的叙事框架。}

框架——选择和强调现实的某些方面以塑造解读的过程~\citep{Entman1993}——是说服的核心机制。叙事使框架具体化且令人难忘。\citet{Green2000}表明，沉浸于故事会降低批判性评估并促进与故事一致的信念改变，这一发现得到了叙事说服更广泛证据的支持~\citep{Shen2015}。在在线社交媒体中，此类策略已成为可规模化的意见操纵工具~\citep{Tucker2018,Wardle2017}。不同叙事依赖不同机制：受害者故事引发共情认同~\citep{Green2000}，阴谋论在不确定性下提供连贯的因果解释~\citep{Douglas2017}，恐惧诉求提升威胁感知~\citep{Witte1992}，而基于权威的主张利用可信度启发式~\citep{Cialdini2001}。近期NLP工作通过宣传技术检测~\citep{DaSanMartino2019}、多维宣传分析~\citep{liu-etal-2025-propainsight}、政治话语的叙事框架模型~\citep{otmakhova-frermann-2025-narrative}，以及跨语言的极端党派和阴谋叙事大规模资源~\citep{maggini-etal-2026-partisanlens}操作化了这些类别。然而，这些工作大多将操纵作为文本分析进行回顾性研究。对于叙事策略一旦进入重复多方讨论后的差异效果知之甚少。

\paragraph{在线操纵的计算方法。}

在线操纵的计算研究主要遵循检测和建模范式。检测工作识别宣传技术~\citep{DaSanMartino2019}，挖掘虚假信息叙事~\citep{sosnowski-etal-2025-dinam}，将免疫理论整合到基于LLM的虚假信息检测中~\citep{modzelewski-etal-2026-malicious}，并刻画机器人驱动或协调性运动~\citep{Ferrara2016,Woolley2018}。建模工作，尤其是基于智能体的意见动力学~\citep{Deffuant2000,Flache2017}，捕捉社会影响下的趋同和极化，但通常将意见视为标量状态并省略影响发生所通过的语言。与此同时，LLM已被证明是有说服力的传播者：LLM生成的消息可与人类撰写的论证一样有效地改变政策态度~\citep{Bai2025}，GPT-4在个性化时可胜过人类辩论者~\citep{Salvi2025}，模型说服力似乎随能力扩展~\citep{Durmus2024,pauli-etal-2025-measuring}，而人口统计感知的人机LLM交互可驱动回音室和极化模式~\citep{borah-etal-2026-persuasion}。这些研究确立了LLM介导说服的风险，但大多检验单消息或二元交互。我们则研究网络化讨论过程中的操纵，其中消息累积、流通并重塑后续回应。

\paragraph{基于LLM的社会模拟。}

基于LLM的社会模拟提供了研究集体行为的受控方式。\citet{Park2023}为可信的生成式智能体引入了记忆、反思和规划，\citet{Gao2023}将LLM智能体扩展到社交网络模拟。近期研究使用此类智能体考察辩论偏差~\citep{taubenfeld-etal-2024-systematic}、群体从众~\citep{choi-etal-2025-empirical}、回音室~\citep{wang-etal-2025-decoding}以及立场--话题共演化~\citep{zhang-etal-2025-spark}。其他研究关注内容传播与审核~\citep{liu-etal-2025-mosaic}、大规模集体行为~\citep{hashemi-macy-2026-empirical}、新闻失真~\citep{liu-etal-2025-stepwise}和平台级操纵~\citep{Touzel2024}。近期批判性工作还强调此类模拟的有效性取决于对生成式智能体如何真实地再现在线传播进行经验基准测试~\citep{munker-etal-2026-dont}，以及保留开放式自由文本输出而非将意见折叠为多选题探针~\citep{ma-etal-2026-open}。这些工作表明LLM智能体可支持对社会动态（包括错误信息相关场景）的经验探究。仍待探索的是不同叙事攻击策略的比较效果、它们与群体特征的交互，以及对干预的响应。NarrateSim是专为此问题构建的社会模拟框架，用作跨话题、群体、网络位置和干预条件的叙事攻击经验测试平台。
